\chapter{Contexte et mission}
\label{ch:contexte}

\questionchapitre{Dans quel dispositif ces travaux s'inscrivent-ils, et quelle mission leur
avait été confiée ?}

\section{Le programme JUNON}

\textsc{Junon} est l'un des programmes \emph{Ambition Recherche et Développement} financés par
la Région Centre-Val de Loire. Le BRGM le coordonne. Sur cinq ans, il mobilise plus de
\SI{12,3}{\mega\euro}, dont \SI{6}{\mega\euro} apportés par la Région, et réunit dix-neuf
partenaires : universités d'Orléans et de Tours, CNRS, INRAE, INSA Centre-Val de Loire, pôle de
compétitivité DREAM Eau et Milieux, collectivités et entreprises du
territoire~\cite{brgm_junon_programme,brgm_junon_lancement}.

Son objectif est de construire des jumeaux numériques des milieux naturels continentaux, l'eau,
le sol et l'air, et les services qui vont avec. Le programme se décline en projets
collaboratifs. Trois portent sur un compartiment : \textsc{Juneau} pour les eaux souterraines,
\textsc{PrevizO} pour les eaux de surface, \textsc{Januss} pour l'air et les zones humides.
Trois autres fournissent les moyens communs : \textsc{Data} pour la gestion des observations,
\textsc{Prediction} pour la modélisation, et le volet qui assemble les jumeaux.

Le LIFAT participe à ces trois volets transversaux sur 2022--2026~\cite{lifat_junon}. Le volet
\textsc{Prediction}, dont Nicolas Labroche a la responsabilité, traite la prédiction sur séries
temporelles environnementales. C'est là que s'inscrit ce postdoctorat.

\section{Le sujet, et ce qu'il supposait}

Le sujet portait sur l'explicabilité des modèles de prévision appliqués aux nappes. La question
est bien posée. Un réseau de neurones qui prévoit un niveau piézométrique produit un nombre, et
ce nombre ne vaut rien tant qu'on ne sait pas répondre à trois questions que tout hydrogéologue
posera : sur quoi le modèle s'appuie-t-il, que se passerait-il dans d'autres conditions, et
jusqu'où peut-on lui faire confiance en dehors de ce qu'il a vu.

Ces questions ne sont pas théoriques. Les niveaux de nappe déclenchent des arrêtés de
restriction, conditionnent des autorisations de prélèvement, dimensionnent des ouvrages
d'alimentation en eau potable. Une prévision qu'on ne sait pas justifier ne sera pas utilisée, et
elle ne devrait pas l'être.

Mais l'explicabilité s'applique à un modèle, et un modèle s'entraîne sur des données. Au
démarrage des travaux, aucun jeu de données n'associait les chroniques piézométriques françaises
à leur forçage climatique. Le chapitre~\ref{ch:verrou} décrit ce manque en détail. Il a
déterminé toute la trajectoire de l'année.

\section{Une mission d'ingénierie, assumée et actée}

Le travail a donc pris une forme que le sujet initial ne laissait pas prévoir : identifier et
qualifier les sources, construire l'entrepôt, produire et valider les indicateurs, puis
seulement outiller la modélisation et l'explicabilité.

Ce n'est pas une reconstruction après coup. La réunion de cadrage du 13 février 2026 entre le
BRGM et l'université de Tours inscrit parmi ses prochaines étapes la « mise en place d'une stack
pérenne pour stocker et harmoniser les données communes », et en attribue nommément le pilotage
à l'auteur de ce rapport~\cite{reunion_aida}.

L'ordre choisi n'est pas un détour. C'est le seul dans lequel les résultats obtenus en bout de
chaîne restent comparables et rejouables. Un modèle entraîné sur une extraction faite pour
l'occasion produit un chiffre que personne ne pourra confronter à un autre.

\begin{remarque}[title={Le fil de ce rapport}]
Le document suit l'ordre dans lequel les obstacles se sont présentés, de la donnée dispersée
jusqu'à la prévision explicable, et non l'ordre des dépôts logiciels. Chaque chapitre s'ouvre
sur la question qu'il traite et se referme sur ce qui est acquis et ce qui reste ouvert.
\end{remarque}

\section{Périmètre}

Deux réalisations logicielles sont couvertes, développées intégralement pendant la période.

\begin{center}
\begin{tabularx}{\textwidth}{@{}lX@{}}
\toprule
\textbf{Dépôt} & \textbf{Objet} \\
\midrule
\hubeau{}~\cite{depot_entrepot} & L'entrepôt : ingestion, transformation et exposition des
données piézométriques, hydrométriques, climatiques et hydrogéologiques. Premier commit le
23 septembre 2025. \\
\junonrepo{}~\cite{depot_plateforme} & La plateforme : observatoire cartographique, laboratoire
de modèles métier, laboratoire d'apprentissage, explicabilité. Premier commit le 6 novembre
2025. \\
\bottomrule
\end{tabularx}
\end{center}

Les deux dépôts totalisent environ \num{1400} contributions. Le contrat s'est achevé le 30 juin
2026. La documentation, la neutralisation des éléments sensibles et la passation se sont
poursuivies jusqu'au 24 août 2026, pour que ces réalisations restent exploitables une fois leur
auteur parti.

\begin{acquisouvert}
\textbf{Acquis.} La mission est claire et sourcée : construire l'infrastructure de données et
d'expérimentation qui manquait, avant de pouvoir traiter le sujet d'explicabilité.

\textbf{Ouvert.} Les travaux d'explicabilité eux-mêmes se poursuivront au-delà de cette période.
Le chapitre~\ref{ch:socle} précise ce qui leur est transmis.
\end{acquisouvert}
